% =====================================================================
%  9. Failure modes and pathological cases
% =====================================================================
\section{When gradient descent goes wrong}

\begin{frame}{Blow-ups, and functions with no global $L$}
  \begin{columns}[T]
    \begin{column}{0.5\textwidth}
      \textbf{1. Step too large} ($\alpha>2/L$): oscillating blow-up along the top eigenvector (M3(f), A3(d)).

      \vspace{4pt}
      \textbf{2. No global $L$.} For $f(x) = x^4$, $f'' = 12x^2$ keeps growing. With $\alpha = 0.1$,
      \[
        x_{k+1} = x_k\,(1 - 0.4x_k^2),
      \]
      which diverges exactly when $|x_0|>\sqrt5\approx2.236$:

      \vspace{2pt}
      \footnotesize
      \begin{tabular}{@{}l l@{}}
        $x_0 = 2.2$: & $-2.059,\ 1.434,\ 0.255,\ 0.249,\dots\to0$\\
        $x_0 = 2.3$: & $-2.567,\ 4.198,\ -25.39,\ 6521,\dots$
      \end{tabular}
    \end{column}
    \begin{column}{0.48\textwidth}
      \textbf{3. A flat minimum.} At $x^\star = 0$ we have $f''(0) = 0$, so $f$ is not strongly convex and the rate
      drops to \emph{sublinear}:
      \[
        x_{k+1}\approx x_k - 0.4x_k^3 \;\Rightarrow\; x_k\approx\frac{1}{\sqrt{0.8\,k}} .
      \]
      \footnotesize Check: at $k = 10^4$ the iteration gives $0.011176$ and the formula $0.011180$.

      \vspace{4pt}
      \begin{warnbox}[Lesson]
        \small A safe step depends on \emph{where you are}. A fixed $\alpha$ is only safe with a global $L$;
        otherwise let a line search pick the step (Section~10).
      \end{warnbox}
    \end{column}
  \end{columns}
\end{frame}

\begin{frame}{Saddle points: a trap with a trapdoor}
  \begin{columns}[T]
    \begin{column}{0.5\textwidth}
      \small
      $f = x^4+y^4-4xy$ has a saddle at $\mathbf 0$, with Hessian eigenvalues $-4$ along $(1,1)$ and $+4$ along $(1,-1)$.
      \begin{itemize}\footnotesize
        \item On the \textbf{stable manifold} $y=-x$, GD slides \emph{into the saddle} (M1).
        \item Elsewhere the unstable part grows by $1+4\alpha$ per step: escaping from distance $\delta = 10^{-3}$ takes
              about 21 steps ($\alpha = 0.1$).
      \end{itemize}
      \begin{thmbox}[Theorem \normalfont\cite{lee2016}]
        \footnotesize For $C^2$ $f$ with $L$-Lipschitz gradient and $\alpha<1/L$, GD from a random start
        converges to a strict saddle with probability zero.
      \end{thmbox}
    \end{column}
    \begin{column}{0.48\textwidth}
      \centering
      \includegraphics[width=0.7\linewidth]{fig_landscape_surface}

      \footnotesize\textcolor{mGrey}{GD almost never \emph{ends} at a saddle, but can idle near one: $\norm{\nabla f}$ is tiny there.}
    \end{column}
  \end{columns}
\end{frame}

\begin{frame}{Local minima: the starting point decides the answer}
  \begin{columns}[T]
    \begin{column}{0.52\textwidth}
      \centering
      \includegraphics[width=\linewidth, height=0.6\textheight, keepaspectratio]{fig_local_minima}
    \end{column}
    \begin{column}{0.46\textwidth}
      \small
      $f(x) = x^4 - 6x^3 + 8x^2$ has two valleys with a hill at $x\approx1.22$.
      \begin{itemize}\footnotesize
        \item From $x_0 = 1.0$ GD rolls left into the \emph{local} minimum $x = 0$.
        \item From $x_0 = 1.5$ it rolls right into the \emph{global} minimum $x\approx3.28$.
      \end{itemize}
      GD may find \emph{a} minimum and miss \emph{the} best one.
      \begin{warnbox}[In practice]
        \footnotesize Try several starts and keep the best. The SSR of a line is convex: one minimum only.
      \end{warnbox}
    \end{column}
  \end{columns}
\end{frame}

\begin{frame}{Plateaus and curved valleys}
  \begin{columns}[T]
    \begin{column}{0.4\textwidth}
      \small
      \textbf{Plateau} (A4). For large $k$ the model $e^{-kt}$ is already close to zero, so it barely reacts to $k$
      and $J'\approx 10^{-3}$. A thousand steps from $k_0 = 6$ only reach $k\approx4.65$: a \alert{vanishing gradient}.

      \vspace{4pt}
      \textbf{The banana valley} (Rosenbrock)
      \[
        f = (1-x)^2+100\,(y-x^2)^2
      \]
      has $\kappa\approx2508$ at the minimum, and the valley \emph{bends}, so no single fixed rescaling straightens it.
      On the right: GD with Armijo against momentum over 2000 iterations.
    \end{column}
    \begin{column}{0.58\textwidth}
      \centering
      \includegraphics[width=\linewidth]{fig_rosenbrock}

      \footnotesize\textcolor{mGrey}{The A4 plateau itself is plotted in Problem A4 (solution 3 of 4).}
    \end{column}
  \end{columns}
\end{frame}

\begin{frame}{Nonsmooth objectives: a constant step never settles}
  $f(x) = |x|$ has ``gradient'' $\operatorname{sign}(x)$, whose size never shrinks near the minimum:
  $\;x_{k+1} = x_k - \alpha\operatorname{sign}(x_k)$.
  \vspace{4pt}
  \begin{itemize}\small
    \item The iterates march into the band $|x|\le\alpha$, then \alert{bounce around in it forever}: the error stays about $\alpha$.
    \item In the smooth case $\nabla f\to\mathbf 0$ shrinks the steps for free (Section~4). Here it does not.
    \item Remedy: shrink the steps yourself (the subgradient method), at the slower rate $O(1/\sqrt k)$~\cite{nesterov2018,boyd2004}.
  \end{itemize}
  \begin{keybox}[Where you meet this]
    \small $\ell_1$ regularisation, hinge loss, ReLU networks, max-type objectives. Smoothness is what makes a constant step work.
  \end{keybox}
\end{frame}

\begin{frame}{Finite precision: gradients by finite differences}
  \begin{columns}[T]
    \begin{column}{0.5\textwidth}
      \small
      No gradient formula? Try $\;f'(x)\approx\frac{f(x+h)-f(x)}{h}$, with error
      \[
        E(h)\approx\underbrace{\tfrac h2|f''|}_{\text{truncation}}+\underbrace{\tfrac{2\epsilon_{\text{mach}}|f|}{h}}_{\text{round-off}},
      \]
      smallest at $h^\star\sim\sqrt{\epsilon_{\text{mach}}}\approx1.5\times10^{-8}$. Going smaller makes it \emph{worse}.
      \begin{itemize}\footnotesize
        \item Forward: about 8 digits at best. Central: $h^\star\sim\epsilon^{1/3}$, error $\sim\epsilon^{2/3}\approx4\times10^{-11}$.
        \item Cost: $n$ extra $f$-evaluations per gradient (autodiff: a constant factor, Section~13).
        \item $\norm{\nabla f}$ has a noise floor: use relative stopping tests.
      \end{itemize}
    \end{column}
    \begin{column}{0.48\textwidth}
      \centering
      \includegraphics[width=\linewidth]{fig_fd_error}

      \footnotesize The classic V: truncation error on the right, round-off on the left ($f = e^x$ at $x = 1$).
    \end{column}
  \end{columns}
\end{frame}
